\scanpage{65}
\mspage{32}
Chapter II. \uline{Sheaves}

\S\ 1. \uline{Existence Theorems}: Let $X$ be a set. The collection of all subsets of $X$ can be considered as the collection of objects of a category $\underline X$ where there is a unique morphism $A\longrightarrow B$ if and only if $B\supset A$. (Note the reversal of order.) If $T_X$ is a topology on $X$ (i.e., a collection of open sets) then $T_X$ determines a full subcategory $\underline T_X$ (or just $\underline T$ when only one space is being considered) of $\underline X$. $\underline X$ and $\underline T$ are clearly small categories. $\underline X$ is complete, for a left (resp., right) root in $\underline X$ is simply a union (resp., intersection), while $\underline T$ is left complete and finitely right complete.

Let $\underline U=\{U_i\}$ be a collection of open sets with $U=\bigcup U_i$. If $\underline U$ is closed under finite intersections (i.e., the subcategory of $\underline T$ determined by $\underline U$ is finitely right complete), then $\underline U$ is called a \uline{strong covering} of $U$. Note that if $D_{\underline U}:\underline U\longrightarrow\underline T$ is the inclusion functor then $U=\mathrm{Lr}\,D_{\underline U}$.

1.1 \uline{Definition}: For any category $\underline A$, the functor category $\underline F(\underline T_X,\underline A)$ will be called the category of \uline{presheaves on} $X$ \uline{with values in} $\underline A$. It will be denoted by $\underline P(X,\underline A)$. The full
\mspage{64}
